
The analysis presented in this dissertation was performed using data taken by \ac{sPHENIX} detector. The detector was located within the \ac{RHIC} complex at Brookhaven National Laboratory. It was installed at 8~o\textquotesingle clock interaction point of the \ac{RHIC} ring, replacing its predecessor ``\ac{PHENIX}'' detector \cite{obrien2024sphenix, belmont2024sphenix}. This chapter begins with an overview of the \ac{RHIC} complex, particle acceleration (for both protons and heavy ions) proceeds with a summary of \ac{sPHENIX} detector, and ends with a detailed explanation of \ac{EMCal} subsystem.


\section{The RHIC Complex}

The \ac{RHIC} accelerator was built on the legacy of the \ac{AGS}, later of which was operational since 1960s \cite{harrison2003rhic, NLRP_2015}. The \ac{AGS} revolutionized the concept of particle acceleration by using the principle of strong focusing (or Alternate Gradient). It utilized an alternating sequence of quadrupole electromagnets to constrain the transverse spread of the particle beam and achieve the unprecedented energy in particle acceleration \cite{aps_bnl_hist}.

\section{Proton and Heavy-Ion Acceleration at RHIC}

Although both protons and heavy ions (such as $Au$) utilize same \ac{RHIC} accelerator facility, their acceleration processes involve fundamentally different approaches. Heavy-ion acceleration focuses on progressive charge stripping, control of space-charge effects and management of intrabeam scattering and emittance growth \cite{harrison2003rhic}. In contrast, proton acceleration focuses on preserving and manipulating the beam polarization for spin-physics experiments \cite{harrison2003rhic}. The proton acceleration requires specialized polarized ion sources, spin rotators, and Siberian snake magnets (in \ac{AGS} and \ac{RHIC} rings) to overcome depolarizing resonances \cite{harrison2003rhic}.


